%% sections/Method.tex
\section{Research Methodology}
\label{sec:method}

\section{\texorpdfstring{\hl{4.
Methodology}}{4. Methodology}}\label{methodology}

\subsection{\texorpdfstring{\hl{4.1 Overall
Workflow}}{4.1 Overall Workflow}}\label{overall-workflow}

\hl{The methodological workflow consists of six stages (As shown in Fig.
6): UAV image acquisition, empirical corpus construction, camera pose
estimation and sparse reconstruction, neural rendering/reconstruction
using four methods, quantitative and qualitative evaluation, and
conservation-oriented interpretation. This workflow follows an
engineering-informatics logic: a real heritage corpus is constructed
first, multiple computational methods are tested under comparable
conditions, and the results are then interpreted in relation to digital
conservation tasks.}

\begin{figure}
\centering
\pandocbounded{\includegraphics[keepaspectratio]{figures/media/image8.emf}}
\caption{\textbf{\hl{Fig. 6. Overall methodological workflow for
UAV-based neural reconstruction and conservation-oriented method
selection.}}}
\end{figure}

\subsubsection{\texorpdfstring{\textbf{\hl{4.2 Benchmark protocol and
fairness
control}}}{4.2 Benchmark protocol and fairness control}}\label{benchmark-protocol-and-fairness-control}

\hl{To ensure that the comparison among NeRF, 2DGS, 3DGS, and INGP was
as fair and reproducible as possible, all methods were evaluated under a
unified benchmark protocol. For each drum tower scene, the same UAV
image set was used as input after image quality screening and resolution
standardization. All images were resized to 1920 × 1080 pixels, and
camera poses were estimated using the same SfM-based workflow to avoid
method-specific pose estimation bias. The resulting camera parameters
and image lists were then converted into the required input formats for
the four reconstruction methods.}

\hl{The train/test split followed an LLFF-style strategy where
applicable, in which one image was selected as a test view at regular
intervals. The same training and test image lists were used for all
methods within each scene. For the Chaoli scene, no independent test
image was available under the current split; therefore, the reported
quantitative values for this scene are interpreted as reference values
rather than fully independent test-set performance. This limitation is
explicitly considered in the results and discussion sections.}

\hl{To reduce implementation-related bias, all methods were trained and
evaluated on the same workstation using identical image resolution,
identical scene-level input data, and comparable camera poses. Default
or officially recommended parameter settings were adopted unless a
method failed to converge under the default configuration. Rendering
quality was evaluated on the same held-out views for each scene where
test views were available. PSNR, SSIM, and LPIPS were computed between
the rendered images and the corresponding reference images. Training
time and rendering speed were recorded under the same hardware
environment to support efficiency comparison.}

\hl{It should be noted that the benchmark is primarily designed to
evaluate image-level rendering quality and practical reconstruction
efficiency for conservation-oriented visualization. Therefore, the
results should not be interpreted as a comprehensive geometric accuracy
ranking of all methods. In particular, 2DGS has a surface-oriented
formulation, and its potential advantages may be more fully reflected
through geometry-specific indicators such as surface consistency, mesh
accuracy, or component-level geometric completeness. This issue is
further discussed as a limitation and as a direction for future work.}

\subsection{\texorpdfstring{\hl{4.3Evaluation
Metrics}}{4.3Evaluation Metrics}}\label{evaluation-metrics}

\hl{Three image-quality metrics were used to evaluate rendered outputs
against reference images: PSNR, SSIM, and LPIPS. PSNR measures
pixel-level fidelity and is reported in dB, where higher values indicate
lower reconstruction error. SSIM measures structural similarity in the
range of 0 to 1, where higher values indicate better structural
consistency. LPIPS measures perceptual feature distance, where lower
values indicate better perceptual similarity. These metrics were
complemented by training time, GPU memory consumption, rendering speed,
and component-level visual interpretation.}

\begin{longtable}[]{@{}
  >{\centering\arraybackslash}p{(\linewidth - 6\tabcolsep) * \real{0.1553}}
  >{\centering\arraybackslash}p{(\linewidth - 6\tabcolsep) * \real{0.3590}}
  >{\centering\arraybackslash}p{(\linewidth - 6\tabcolsep) * \real{0.1462}}
  >{\centering\arraybackslash}p{(\linewidth - 6\tabcolsep) * \real{0.3396}}@{}}
\caption{\textbf{\hl{Table 7. Evaluation metrics and
interpretation}}}\tabularnewline
\toprule\noalign{}
\begin{minipage}[b]{\linewidth}\centering
\textbf{\hl{Metric}}
\end{minipage} & \begin{minipage}[b]{\linewidth}\centering
\textbf{\hl{Full name}}
\end{minipage} & \begin{minipage}[b]{\linewidth}\centering
\textbf{\hl{Direction}}
\end{minipage} & \begin{minipage}[b]{\linewidth}\centering
\textbf{\hl{Interpretation}}
\end{minipage} \\
\midrule\noalign{}
\endfirsthead
\toprule\noalign{}
\begin{minipage}[b]{\linewidth}\centering
\textbf{\hl{Metric}}
\end{minipage} & \begin{minipage}[b]{\linewidth}\centering
\textbf{\hl{Full name}}
\end{minipage} & \begin{minipage}[b]{\linewidth}\centering
\textbf{\hl{Direction}}
\end{minipage} & \begin{minipage}[b]{\linewidth}\centering
\textbf{\hl{Interpretation}}
\end{minipage} \\
\midrule\noalign{}
\endhead
\bottomrule\noalign{}
\endlastfoot
\hl{PSNR} & \hl{Peak Signal-to-Noise Ratio} & \hl{Higher is better} &
\hl{Pixel-level reconstruction fidelity} \\
\hl{SSIM} & \hl{Structural Similarity Index Measure} & \hl{Higher is
better} & \hl{Structural and visual consistency} \\
\hl{LPIPS} & \hl{Learned Perceptual Image Patch Similarity} & \hl{Lower
is better} & \hl{Perceptual similarity and texture realism} \\
\hl{Training time} & \hl{Average training time} & \hl{Lower is better} &
\hl{Computational efficiency} \\
\hl{Rendering speed} & \hl{Frames per second} & \hl{Higher is better} &
\hl{Real-time visualization capability} \\
\end{longtable}

\subsection{\texorpdfstring{\hl{4.4 Experimental Setup and
Implementation
Details}}{4.4 Experimental Setup and Implementation Details}}\label{experimental-setup-and-implementation-details}

\hl{All experiments were conducted on a GPU cloud server equipped with
an NVIDIA GeForce RTX 4090 GPU with 24 GB memory. Four neural
rendering/reconstruction methods were compared: NeRF, 2DGS, 3DGS, and
INGP. All methods used the same scene-level image corpus and the same
train/test split scheme wherever applicable, and the final reported
metrics were averaged over all available evaluation images.}

\hl{Camera pose estimation and sparse reconstruction were performed
using COLMAP 3.14.0 with CUDA acceleration. SIFT features were extracted
with GPU acceleration, and the Sequential Matcher was adopted because
the input images were derived from video sequences or sequential UAV
acquisition. The OpenCV camera model was used, including the intrinsic
and distortion parameters fx, fy, cx, cy, k1, k2, p1, and p2. A
single-camera constraint was enabled using
-\/-ImageReader.single\_camera 1. Sparse reconstruction was carried out
using the COLMAP incremental SfM mapper, and image undistortion was
performed using COLMAP image\_undistorter.}

\begin{longtable}[]{@{}
  >{\centering\arraybackslash}p{(\linewidth - 8\tabcolsep) * \real{0.0660}}
  >{\centering\arraybackslash}p{(\linewidth - 8\tabcolsep) * \real{0.2140}}
  >{\centering\arraybackslash}p{(\linewidth - 8\tabcolsep) * \real{0.3892}}
  >{\centering\arraybackslash}p{(\linewidth - 8\tabcolsep) * \real{0.1531}}
  >{\centering\arraybackslash}p{(\linewidth - 8\tabcolsep) * \real{0.1777}}@{}}
\caption{\textbf{\hl{Table8. Reconstruction methods and implementation
settings}}}\tabularnewline
\toprule\noalign{}
\begin{minipage}[b]{\linewidth}\centering
\textbf{\hl{Method}}
\end{minipage} & \begin{minipage}[b]{\linewidth}\centering
\textbf{\hl{Framework / codebase}}
\end{minipage} & \begin{minipage}[b]{\linewidth}\centering
\textbf{\hl{Key settings}}
\end{minipage} & \begin{minipage}[b]{\linewidth}\centering
\textbf{\hl{Output resolution}}
\end{minipage} & \begin{minipage}[b]{\linewidth}\centering
\textbf{\hl{Protocol}}
\end{minipage} \\
\midrule\noalign{}
\endfirsthead
\toprule\noalign{}
\begin{minipage}[b]{\linewidth}\centering
\textbf{\hl{Method}}
\end{minipage} & \begin{minipage}[b]{\linewidth}\centering
\textbf{\hl{Framework / codebase}}
\end{minipage} & \begin{minipage}[b]{\linewidth}\centering
\textbf{\hl{Key settings}}
\end{minipage} & \begin{minipage}[b]{\linewidth}\centering
\textbf{\hl{Output resolution}}
\end{minipage} & \begin{minipage}[b]{\linewidth}\centering
\textbf{\hl{Protocol}}
\end{minipage} \\
\midrule\noalign{}
\endhead
\bottomrule\noalign{}
\endlastfoot
\hl{NeRF} & \hl{TensorFlow 2.x; modified nerf-pytorch / bmild-style} &
\hl{200,000 iterations; exponential learning-rate decay; positional
encoding with 10 bands; checkpoint resumption} & \hl{Original setting /
{[}TO BE FILLED{]}} & \hl{Final checkpoint rendered on evaluation
set} \\
\hl{2DGS} & \hl{Official PyTorch; hbb1/2d-gaussian-splatting} &
\hl{30,000 iterations; depth-loss weight; SfM sparse point cloud
constraint} & \hl{Half resolution (-\/-resolution 2)} & \hl{Final
checkpoint rendered on evaluation set} \\
\hl{3DGS} & \hl{Official PyTorch; graphdeco-inria/gaussian-splatting} &
\hl{30,000 iterations; default anisotropic regularization} & \hl{Half
resolution (-\/-resolution 2)} & \hl{Final checkpoint rendered on
evaluation set} \\
\hl{INGP} & \hl{Instant-NGP v2.0dev} & \hl{Multiresolution hash
encoding; 16 logarithmic levels} & \hl{{[}TO BE FILLED{]}} &
\hl{Built-in rendering script for evaluation} \\
\end{longtable}

\begin{longtable}[]{@{}
  >{\centering\arraybackslash}p{(\linewidth - 2\tabcolsep) * \real{0.2618}}
  >{\centering\arraybackslash}p{(\linewidth - 2\tabcolsep) * \real{0.7382}}@{}}
\caption{\textbf{\hl{Table9. Hardware, software, and camera-pose
estimation environment}}}\tabularnewline
\toprule\noalign{}
\begin{minipage}[b]{\linewidth}\centering
\textbf{\hl{Item}}
\end{minipage} & \begin{minipage}[b]{\linewidth}\centering
\textbf{\hl{Configuration}}
\end{minipage} \\
\midrule\noalign{}
\endfirsthead
\toprule\noalign{}
\begin{minipage}[b]{\linewidth}\centering
\textbf{\hl{Item}}
\end{minipage} & \begin{minipage}[b]{\linewidth}\centering
\textbf{\hl{Configuration}}
\end{minipage} \\
\midrule\noalign{}
\endhead
\bottomrule\noalign{}
\endlastfoot
\hl{GPU} & \hl{NVIDIA GeForce RTX 4090} \\
\hl{GPU memory} & \hl{24 GB} \\
\hl{Frameworks} & \hl{NeRF: TensorFlow 2.x; 2DGS/3DGS: PyTorch; INGP:
Instant-NGP CUDA} \\
\hl{Camera pose estimation} & \hl{COLMAP 3.14.0 with CUDA} \\
\hl{Feature extraction} & \hl{SIFT with GPU acceleration} \\
\hl{Matching strategy} & \hl{Sequential Matcher} \\
\hl{Camera model} & \hl{OpenCV (fx, fy, cx, cy, k1, k2, p1, p2)} \\
\hl{Single-camera constraint} & \hl{Enabled
(-\/-ImageReader.single\_camera 1)} \\
\hl{Sparse reconstruction} & \hl{Incremental SfM} \\
\hl{Image undistortion} & \hl{COLMAP image\_undistorter} \\
\end{longtable}
